\section{Introduction}

Text embedding models serve as the fundamental backbone for a wide array of AI applications, including semantic search, retrieval-augmented generation (RAG), text classification, and clustering. By mapping unstructured text into dense vector spaces, these models allow machines to capture complex semantic relationships, enabling efficient and accurate information retrieval and data analysis across massive datasets. This field has recently transitioned from encoder-based architectures~\citep{2019BERT,2019RoBERTa,2020XLM-R} to decoder-based LLM embeddings~\citep{2025Qwen3-Embedding,2025NV-Embed,2025f2llm}, benefiting from the extensive reasoning and linguistic capabilities acquired during large-scale pre-training and achieving remarkable gains in performance.

Despite these advancements, the current state of frontier embedding research is characterized by two significant limitations. First, there is a pervasive English-centric bias in both model training and benchmark evaluation. While benchmarks such as MTEB have been instrumental in standardizing evaluation, the high-resource language subsets therein - such as English and Chinese - receive a disproportionately large share of attention, resulting in an abundance of models that are performant in English but fail to provide global utility. Second, a transparency gap has emerged within the research community. Most top-performing embedding models, such as Gemini-Embedding~\citep{2025Gemini-Embedding} and Qwen3-Embedding~\citep{2025Qwen3-Embedding}, are released either as closed-source APIs or open-weight models without disclosing the underlying training data or methodologies. This lack of transparency hinders reproducibility and limits our collective understanding of how to build truly inclusive, general-purpose embedding systems.

To directly tackle these challenges, we introduce F2LLM-v2, a new family of general-purpose, multilingual embedding models designed to address these critical imbalances. We curate a massive, high-quality training corpus of 60 million samples spanning 282 natural languages and over 40 programming languages solely from publicly available resources. By prioritizing real-world data availability over benchmark-specific optimization, we create a model family that excels across a truly global range of applications, including those involving underserved languages. Besides linguistic inclusivity, we also address computational inclusivity by providing 8 distinct model sizes, ranging from 80M to 14B parameters. By integrating Matryoshka Representation Learning (MRL) and a two-stage training pipeline enhanced by model pruning and novel knowledge distillation, we ensure high performance even in resource-constrained environments. Extensive evaluations confirm that our 14B model achieves state-of-the-art results on 11 MTEB benchmarks, setting a new standard for multilingual embedding capabilities, while the smaller models also outperform previous frontier models with a similar size. To foster an open and equitable research environment, we release the complete training recipe, intermediate checkpoints, and all associated code and data for the F2LLM-v2 family, aiming to drive progress toward a more inclusive future for AI technology.